\begin{textsample}{POS Dim 1 – pi – Score 38.00 – pi/pi\_ef\_36.txt}  \label{ex:f1_pos_088}
Conceito Concepções \textbf{teóricas} Autores/referências¹ Lugar Muitos \textbf{autores} utilizam o termo lugar para se referir à ideia de pertencimento ( TUAN, 1983 ; SCARLATO, 2005 ; \textbf{OLIVEIRA}, 2000, 2013 ; FURLAN, 2004 ).
Lugar seria a expressão do espaço \textbf{vivido}, percebido e representado.
Nessa \textbf{abordagem}, lugar ganha sentido de leitura perceptiva e de campo simbólico. \textbf{Uma} pessoa \textbf{vive} num local, mas o lugar seria sua identificação afetiva, a ligação e vínculo com a paisagem.
Para Edward Relph ( 2012 ), o lugar é \textbf{um} microcosmo.
É onde cada \textbf{um} de nós se relaciona com o mundo e onde o mundo se relaciona conosco.
Acontece, por exemplo, quando \textbf{um} teresinense acessa \textbf{um} site chinês de compras, e adquire \textbf{um} produto produzido naquele país e o recebe em sua casa.
Autores/referências : - CLAVAL, Paul - CORREA, Roberto L. - FURLAN, Sueli A. - \textbf{OLIVEIRA}, Lívia - \textbf{SANTOS}, Milton - SCARLATO, Francisco - C.
TUAN, Yi-Fu.
Território O conceito de território pode ser definido a partir de distintos pontos de vista, pois a Geografia não tem exclusividade em relação a ele.
Diversas áreas do conhecimento utilizam o conceito de território de acordo com sua própria perspectiva \textbf{predominante}.
Por exemplo, a Ciência Política tende a valorizar a perspectiva ligada às relações de poder, principalmente no \textbf{que} \textbf{diz} respeito aos Estados ; a Antropologia tende a valorizar aspectos ligados à cultura e ao simbolismo dos povos ; a Biologia considera os aspectos naturais ; a Psicologia, as dimensões da construção da identidade do indivíduo.
Na Geografia, território é o produto da materialidade técnica das sociedades.
É também campo de forças políticas onde as ações humanas constroem as \textbf{marcas} de sua produção e projetam sua cultura.
Autores/referências : - ANDRADE, Manuel C. - COSTA, Paulo G. - HAESBAERT, Rogério - MORAES Antonio Carlos R. - \textbf{SANTOS}, Milton Região Conceito \textbf{historicamente} utilizado em Geografia, \textbf{que} inicialmente considerava os atributos naturais como diferenciados dos espaços geográficos.
Corrêa ( 1989 ) considera região \textbf{uma} entidade concreta, resultado de múltiplas determinações.
A região não é \textbf{uma} unidade \textbf{que} contém \textbf{uma} diversidade, mas é produto de \textbf{uma} operação de homogeneização, \textbf{que} se dá na luta com as forças \textbf{que} dominam outros espaços regionais, por isso ela é aberta, móvel e \textbf{atravessada} por diferentes relações de poder ( ALBUQUERQUE JÚNIOR, 1999, p. 24 ).
Autores/referências : - CORRÊA, Roberto L. - LENCIONE, Sandra - HASBAERT, R. - \textbf{SANTOS}, Milton - LA BLACHE, V. - MOREIRA, Ruy - GOMES Paulo C. - RIBEIRO, Luiz A de M. - RUA, João Natureza Cada período histórico é marcado por \textbf{um} determinado \textbf{posicionamento} filosófico em relação à concepção de natureza.
As explicações e as definições de natureza acompanham as concepções de mundo \textbf{dependendo} do grupo humano, do tipo de sociedade ou da classe social de quem responde ( CARVALHO, 1991 ).
A forma de estudar e interpretar os sistemas naturais segue essa ampla gama de construções epistemológicas.
A natureza é \textbf{uma} construção social da interpretação dos sistemas naturais.
Em Geografia, estuda -se tanto os sistemas em si, como as ideias de natureza.
A partir dessa construção humana, estabelecemos formas de concebê -la e de nos relacionarmos com o ambiente.
Na atualidade, evidencia -se em diversas áreas do conhecimento a eclosão de novas \textbf{teorias} ( Teoria da Auto-organização, Teoria da Complexidade, Teoria das Estruturas Dissipativas etc. ) referentes a essas novas visões de mundo \textbf{que} consequentemente trazem consigo novas concepções acerca da natureza.
A Geografia trabalha com \textbf{uma} conceituação ampla de natureza : \textbf{funcional}, simbólica, sagrada e produzida pelo capitalismo.
Autores/referências : - CARVALHO, Marcos B. - HASSLER, Márcio L. - LENOBLE, Robert - MORIN, Edgar - \textbf{SANTOS}, Milton - VITTE, Antonio C. - CIGOLONI, Adilar - SCHELLMANN, Karin - VESENTINI, José W. - WHITEHEAD, Alfred N. ¹ – Além dos \textbf{autores} citados nas concepções \textbf{teóricas}, foram acrescidos outros \textbf{que} contribuem para \textbf{uma} melhor compreensão dos conceitos-chaves de Geografia.
Fonte : organizado pelos \textbf{autores} ( 2018 ) Ao utilizar corretamente esses conceitos, desenvolvendo procedimentos de pesquisa e análise geográficas, os estudantes podem reconhecer, por exemplo, as diversas formas de manifestação das desigualdades sociais e regionais, bem como seus fatores condicionantes e as implicações deles decorrentes.
Desse modo, a aprendizagem da Geografia estimula a capacidade de empregar o \textbf{raciocínio} geográfico para pensar e resolver problemas gerados na vida cotidiana, condição fundamental para o desenvolvimento das competências gerais previstas na BNCC.
É importante ressaltar \textbf{que} o conceito de espaço é inseparável do conceito de tempo e ambos precisam ser pensados articuladamente como \textbf{um} processo.
Assim como para a História, o tempo é para a Geografia \textbf{uma} construção social, \textbf{que} se associa à memória e às identidades sociais dos \textbf{sujeitos}.
Do mesmo modo, os tempos da natureza não podem ser ignorados, pois marcam a memória da Terra, ou seja, as transformações ocorridas e suas implicações nas \textbf{atuais} condições do meio físico natural.
Assim, pensar a temporalidade das ações humanas e das sociedades por meio da relação tempo-espaço representa \textbf{um} importante e desafiador processo na aprendizagem de Geografia.
Para isso, é \textbf{preciso} superar a aprendizagem com base apenas na descrição de informações e fatos do dia a dia, cujo significado \textbf{restringe} -se apenas ao contexto \textbf{imediato} da vida dos \textbf{sujeitos}.
A ultrapassagem dessa condição meramente descritiva \textbf{exige} o domínio de conceitos e generalizações.
Estes permitem novas formas de \textbf{ver} o mundo e de compreender, de maneira ampla e crítica, as múltiplas relações \textbf{que} conformam a realidade, de acordo com o aprendizado do conhecimento da ciência geográfica.
A observação da realidade, própria da vivência, possibilita ao estudante o conhecimento de outros espaços geográficos.
Por isso, a necessidade de desenvolver no aluno a capacidade de \textbf{uma} leitura mais aguçada acerca de diferentes lugares.
Eles devem perceber diferenças nas características físicas e humanas próprias de cada espaço de vivência.
No entanto, os alunos devem ser provocados por meio de diferentes questões \textbf{que} \textbf{exigem} \textbf{uma} mobilização em busca de resposta, \textbf{que} consequentemente \textbf{desencadeia} outras perguntas.
Esse movimento possibilita aos alunos a \textbf{percepção} de \textbf{que} a situação geográfica é fruto de \textbf{um} conjunto de relações.
Desse modo, o estudante precisa estar em contato permanente com o seu objeto do conhecimento por meio dos conteúdos, conceitos e processos da geografia.
A BNCC reitera \textbf{que} o estudante deve ter ciência sobre o \textbf{que} vai aprender e para \textbf{que} serve aquele conteúdo para sua vida diária, assim, encontrará sentido para aquilo \textbf{que} aprende na escola.
O documento determina as aprendizagens essenciais \textbf{que} devem ser asseguradas aos estudantes em todos os níveis da educação básica e estabelece as seguintes Unidades Temáticas no componente Geografia : - O \textbf{sujeito} e seu lugar no mundo - Conexões e escalas - Mundo do trabalho - Formas de representação e pensamento espacial - Natureza, ambientes e qualidade de vida.
As unidades temáticas citadas devem ser articuladas para \textbf{que} se assegurem os direitos de aprendizagem e desenvolvimento, garantindo aos estudantes as condições de aprender e se desenvolver.
Desse modo, em cada \textbf{uma} das Unidades Temáticas ( Anos Iniciais e Finais ) do Ensino Fundamental, as possibilidades de aprendizagem ampliam -se em vários aspectos.
Neste momento da vida escolar, os estudantes já possuem maior autonomia em relação à leitura e à escrita, bem como \textbf{um} maior domínio dos procedimentos de observação, descrição, explicação e representação, fato \textbf{que} permite aos mesmos \textbf{que} construam compreensões mais complexas, permitindo analogias e sínteses mais elaboradas, a partir do desenvolvimento de pesquisas de campo e trabalhos escritos.
Diante do \textbf{exposto}, espera -se \textbf{que} o estudo da Geografia no Ensino Fundamental, em todos os seus anos, possa contribuir para o delineamento do projeto de vida dos estudantes, de modo \textbf{que} eles compreendam a produção social do espaço e a sua transformação. \textbf{Que} entendam também os conceitos-chave da Geografia, com os quais mantêm contato desde os anos inicias de suas vidas escolares, e \textbf{que} os mesmos \textbf{desempenham} importante papel na formação do \textbf{raciocínio} espacial e na prática da cidadania.
Deve -se perceber \textbf{que} no ensino de Geografia há \textbf{uma} \textbf{forte} relação com imagens, gráficos e mapas.
Por isso, o importante é fazer destes elementos, instrumentos \textbf{que} permitam de forma constante \textbf{uma} postura sempre reflexiva e crítica dos estudantes, em relação ao mundo nas suas diferentes escalas.
Para tanto, o componente curricular Geografia está organizado em cinco unidades temáticas, em \textbf{uma} progressão das habilidades.
Na unidade temática O \textbf{sujeito} e seu lugar no mundo, focalizam -se as noções de pertencimento e identidade contextualizadas cultura e espacialmente, ou seja, deve considerar as representações da vida dos estudantes.
Além disso, pretende -se possibilitar \textbf{que} os estudantes construam sua identidade relacionando -se com o outro ; valorizem as suas memórias e \textbf{marcas} do passado vivenciadas em diferentes lugares ; e, à medida \textbf{que} se alfabetizam, ampliem a sua compreensão do mundo.
Em prosseguimento, no Ensino Fundamental, procura -se expandir o olhar para a relação do estudante com contextos mais amplos, considerando temas políticos, econômicos e culturais do Piauí, do Brasil e do mundo.
A partir disso, o estudo da Geografia constitui -se em \textbf{uma} busca do lugar de cada indivíduo no mundo, valorizando a sua individualidade e, ao mesmo tempo, situando -o em \textbf{uma} \textbf{categoria} mais ampla de \textbf{sujeito} social : a de cidadão ativo, democrático e solidário.
Enfim, cidadãos produtos de sociedades localizadas em determinado tempo e espaço, mas também produtores dessas mesmas sociedades com sua cultura e suas regras.
Em Conexões e escalas, a atenção está na articulação de diferentes espaços e escalas de análise, possibilitando \textbf{que} os alunos compreendam as relações existentes entre fatos nos níveis local, regional e global.
Portanto, no decorrer do Ensino Fundamental, os alunos precisam compreender as interações multiescalares existentes entre sua vida familiar, seus grupos e espaços de convivência e as interações espaciais mais complexas.
A conexão é \textbf{um} princípio da Geografia \textbf{que} estimula a compreensão do \textbf{que} ocorre entre os componentes da sociedade e do meio físico natural.
Ela também analisa o \textbf{que} ocorre entre quaisquer elementos \textbf{que} constituem \textbf{um} conjunto na superfície terrestre e \textbf{que} explicam \textbf{um} lugar na sua \textbf{totalidade}.
Conexões e escalas explicam os arranjos das paisagens, a localização e a distribuição de diferentes fenômenos e objetos técnicos, por exemplo.
Na unidade temática Mundo do trabalho, o espaço é resultante das relações de produção e reprodução da vida humana, \textbf{que} se dão através do trabalho.
O mundo, \textbf{que} no início era formado somente por elementos naturais, com a presença do ser humano e os processos por ele engendrados, vai, aos poucos, modificando -se a partir da inserção de objetos técnicos e da substituição do meio natural por \textbf{um} meio cada vez mais artificializado.
Na \textbf{abordagem} das Formas de Representação e Pensamento Espacial, o estudante gradualmente desenvolverá, por meio da linguagem \textbf{cartográfica} e gráfica, o \textbf{raciocínio} geográfico.
O \textbf{que} significa criar condições para \textbf{que} o estudante, inclusive os \textbf{que} possuem deficiência visual, ao tempo em \textbf{que} é levado a aprender sobre a espacialidade dos fenômenos, possa localizá -los, compreendê -los e explicá -los, tornando -se \textbf{um} leitor da realidade \textbf{que} o cerca e do mundo, contribuindo também para \textbf{uma} educação inclusiva, proporcionando maior autonomia, acessibilidade e condições para o exercício da cidadania.
Por fim, na unidade temática Natureza, Ambiente e Qualidade de Vida, no Ensino Fundamental recaem as necessidades de se utilizar a natureza como meio de sobrevivência.
Nesta unidade temática, se destaca a evolução técnica da natureza dos seus primórdios à atualidade.
O estudante pode reconhecer de \textbf{que} forma as diferentes comunidades transformam e valorizam a natureza, tanto em relação às \textbf{inúmeras} possibilidades de uso ao transformá -las em recursos, quantos aos impactos socioambientais delas provenientes.
Dessa forma, é possível compreender a história das relações entre sociedade e natureza através da substituição de \textbf{um} meio natural por \textbf{um} meio cada vez mais técnico-científico-informacional.
Assim, ao se estudarem os objetos do conhecimento de Geografia, a partir das unidades temáticas, a ênfase do aprendizado é na posição relativa e dinâmica dos objetos no espaço e no tempo, o \textbf{que} \textbf{exige} a compreensão das características de \textbf{um} lugar ( localização, extensão, conectividade, entre outras ), resultantes das relações com outros lugares.
Por causa disso, o entendimento da situação geográfica, pela sua natureza, é o procedimento para o estudo dos objetos de aprendizagem pelos estudantes.
Em \textbf{uma} mesma atividade a ser desenvolvida pelo professor, os estudantes podem mobilizar, ao mesmo tempo, diversas habilidades de diferentes unidades temáticas.
Considerando esses pressupostos, e em articulação com as competências gerais da BNCC e com as competências específicas da área de Ciências Humanas, o componente curricular de Geografia também deve garantir aos estudantes o desenvolvimento de competências específicas.

% matched lemmas: abordagem, atravessar, atual, autor, cartográfico, categoria, depender, desempenhar, desencadear, dizer, exigir, exposto, forte, funcional, historicamente, imediato, inúmero, marca, oliveira, percepção, posicionamento, preciso, predominante, que, raciocínio, restringir, santo, sujeito, teoria, teórico, totalidade, um, uma, ver, viver
\end{textsample}
